% Source PDF 313--324 = printed 301--312.
% The source PDF313 begins with the last exercise of Chapter9; preserve its position.
% Source PDF 313 = printed 301.
\begin{origstatement}{习题}{9.6.29}\label{appA:ex9629}
假设 $S$ Noether，$X/S$ 射影且平坦。假设 $\mathbf{Pic}_{X/S}$ 存在并表示 $\Pic_{(X/S)\, (\mathrm{fppf})}$。设 $\Lambda$ 是 $\mathbf{Pic}_{X/S}$ 的任意子集，$\Pi$ 是相应的 Hilbert 多项式集合。证明：若 $\Lambda$ 拟紧，则 $\Pi$ 有限。证明：若 $\Lambda$ 连通，则 $\Pi$ 只有一个元素。
\end{origstatement}

\chapter*{附录 A\quad 全部习题的解答}
\addcontentsline{toc}{chapter}{附录 A 全部习题的解答}
\markboth{附录 A\quad 全部习题的解答}{附录 A\quad 全部习题的解答}

这些习题并非故意设置难点，而是为了帮助你检查、巩固和拓展对这些思想与方法的理解。因此，为了自己的数学成长，请在阅读这里的解答之前，认真尝试完成每道习题。注意，解答的编号与习题的编号相同。

\begin{origstatement}{解答}{9.2.3}\label{appA:ans923}
由于 $A$ 是局部环，$\Pic(T)$ 平凡；因此定义 9.2.1 和 9.2.2 给出第一个同构 $\Pic_X(A)\xrightarrow{\sim}\Pic_{X/S}(A)$。

设 $\mathcal L$ 是 $X_A$ 上的可逆层。假设 $\mathcal L$ 的同构类在 $\Pic_{(X/S)\,(\mathrm{zar})}(A)$ 中映到 $0$。那么存在 Zariski 覆盖 $T'\to T$ 及同构 $v'\colon\mathcal L_{T'}\xrightarrow{\sim}\mathcal O_{X_{T'}}$。$T'$ 是 $T$ 的若干 Zariski 开子概形的不交并。其中一个包含闭点，因而等于 $T$。限制 $v'$，得到同构 $\mathcal L\xrightarrow{\sim}\mathcal O_{X_A}$。所以 $\Pic_X(A)\to\Pic_{(X/S)\,(\mathrm{zar})}(A)$ 是单射。

给定 $\lambda\in\Pic_{(X/S)\,(\mathrm{zar})}(A)$，在适当的 Zariski 覆盖 $T'\to T$ 上，用 $X_{T'}$ 上的可逆层 $\mathcal L'$ 表示 $\lambda$。同样，$T'$ 包含一份作为开闭子概形的 $T$。限制 $\mathcal L'$，得到 $X_A$ 上的可逆层 $\mathcal L$。由于 $\mathcal L'$ 表示 $\lambda$，其到 $T'\times_T T'$ 的两个拉回之间有同构 $v''$。把 $v''$ 限制到 $T'\times_T T$，即 $T'$，得到 $\mathcal L'$ 与 $\mathcal L_{T'}$ 之间的同构。因此 $\mathcal L$ 也表示 $\lambda$。所以 $\Pic_X(A)\to\Pic_{(X/S)\,(\mathrm{zar})}(A)$ 也是满射。

假设 $A$ 是剩余域代数闭的 Artin 局部环。由 [EGAIV4, 17.6.2 和 17.6.3]，每个有限型平展 $A$-代数 $B$ 都是 Artin 局部代数的直积，且每个因子都同构于 $A$。所以，若 $T'\to T$ 是任意平展覆盖，则 $T'$ 是若干开子概形的不交并，每一个都是一份 $T$。因此，按上面的论证，得到 $\Pic_X(A)\xrightarrow{\sim}\Pic_{(X/S)\,(\text{ét})}(A)$。

假设 $A=k$，其中 $k$ 是代数闭域。那么任意 fppf 覆盖 $T'\to T$ 都有截面；事实上，$T'$ 的任意闭点处的局部环在 $k$ 上本质有限型，因此由 Hilbert 零点定理，其剩余域等于 $k$。这个点不必在 $T'$ 中孤立；但作本质上同样的论证，仍得到 $\Pic_X(k)\xrightarrow{\sim}\Pic_{(X/S)\,(\mathrm{fppf})}(k)$。\hfill$\square$
\end{origstatement}

\begin{origstatement}{解答}{9.2.4}\label{appA:ans924}
扩张 $\mathbb C/\mathbb R$ 平展。因此，若 $\varphi^*\mathcal O(1)$ 到 $X_{\mathbb C\otimes_{\mathbb R}\mathbb C}$ 的两个拉回同构，则 $\varphi^*\mathcal O(1)$ 定义 $\Pic_{(X/\mathbb R)\,(\text{ét})}(\mathbb R)$ 的元素 $\lambda$。

取不定元 $z$，把 $\mathbb C$ 与 $\mathbb R[z]/(z^2+1)$ 认同。那么，经纯量扩张，$\mathbb C\otimes_{\mathbb R}\mathbb C$ 与 $\mathbb C[z]/(z-1)(z+1)$ 认同。因此由中国剩余定理，$\mathbb C\otimes_{\mathbb R}\mathbb C$ 同构于积 $\mathbb C\times\mathbb C$。（相应地，$c\otimes1$ 和 $1\otimes c$ 分别与 $(c,c)$ 和 $(c,\bar c)$ 认同，其中 $\bar c$ 是 $c$ 的共轭，但这里不需要这一事实。）
% Source p301 prints (z-1)(z+1), although the declared polynomial is z^2+1; retained.

所以，$X_{\mathbb C\otimes_{\mathbb R}\mathbb C}$ 在 $\mathbb R$ 上同构于两份 $X_{\mathbb C}$ 的不交并。对任意域 $k$，$\mathbb P_k^1$ 上的可逆层 $\mathcal L$ 由一个整数，即其 Euler 示性数 $\chi(\mathcal L)$，唯一确定到同构，见 [Har77, Cor.6.17, p.145]。因此，$\varphi^*\mathcal O(1)$ 到 $X_{\mathbb C\otimes_{\mathbb R}\mathbb C}$ 的两个拉回同构。于是 $\varphi^*\mathcal O(1)$ 定义 $\Pic_{(X/\mathbb R)\,(\text{ét})}(\mathbb R)$ 中的 $\lambda$。
% Source PDF 314 = printed 302.

（类似地，$\varphi^*\mathcal O(1)$ 在 $X_{\mathbb C}$ 上的同构类不依赖同构 $\varphi\colon X_{\mathbb C}\xrightarrow{\sim}\mathbb P_{\mathbb C}^1$ 的选择。因此 $\lambda$ 也不依赖。但这里同样不需要这一事实。）

最后，必须证明 $\lambda$ 不在 $\Pic_{(X/\mathbb R)\,(\mathrm{zar})}(\mathbb R)$ 的像中。反证，假设它在像中。那么 $\lambda$ 来自 $X$ 上的可逆层 $\mathcal L$。事先不能断定拉回 $\mathcal L|X_{\mathbb C}$ 同构于 $\varphi^*\mathcal O(1)$；只要求这两个可逆层在拉回到 $X_A$ 后成为同构，其中 $A$ 是某个平展 $\mathbb C$-代数。

然而，上同调与平坦基变换可交换，所以
\[
 \dim_{\mathbb R}H^0(\mathcal L)
 =\operatorname{rank}_A H^0(\mathcal L|X_A)
 =\dim_{\mathbb C}H^0(\varphi^*\mathcal O(1))=2.
\]
因此 $\mathcal L$ 有非零截面，定义正合序列
\[
 0\longrightarrow\mathcal O_X\longrightarrow\mathcal L\longrightarrow\mathcal O_D\longrightarrow0.
\]
类似地，$H^1(\mathcal O_X)=0$。因此 $\dim_{\mathbb R}H^0(\mathcal O_D)=1$，故 $D$ 是 $X$ 的一个 $\mathbb R$-点。但 $X$ 没有 $\mathbb R$-点。所以 $\lambda$ 不在 $\Pic_{(X/\mathbb R)\,(\mathrm{zar})}(\mathbb R)$ 的像中。\hfill$\square$
\end{origstatement}

\begin{origstatement}{解答}{9.2.6}\label{appA:ans926}
首先，本质上由定义有 $\Pic_{(X/S)\,(\mathrm{fppf})}(k)=\Pic_{(X_k/k)\,(\mathrm{fppf})}(k)$，因为 $k$-概形的映射 $T'\to T$ 是 fppf 覆盖，当且仅当它作为 $S$-概形的映射是 fppf 覆盖。其余三个函子也可作类似分析。$f_k\colon X_k\to k$ 有截面；事实上，$f_k$ 有限型且 $k$ 代数闭，因而由 Hilbert 零点定理，$X$ 的任意闭点的剩余域均为 $k$。于是由定理 9.2.5 的第二部分，四个函子的 $k$-点相同。最后，$\Pic_{X/S}(k)=\Pic(X_k)$，因为只要 $T$ 仅有一个闭点，$\Pic(T)$ 就平凡。

无论 $\mathcal O_S\xrightarrow{\sim}f_*\mathcal O_X$ 是否普遍成立，由习题 9.2.3，四个函子都有相同的几何点；事实上，给定代数闭域 $k$，这些函子的 $k$-点恰好是 $\Pic(X_k)$ 的元素。\hfill$\square$
\end{origstatement}

\begin{origstatement}{解答}{9.3.2}\label{appA:ans932}
由定义，$H^0(X,\mathcal L)_{\mathrm{reg}}$ 中的截面对应于单射 $\mathcal L^{-1}\hookrightarrow\mathcal O_X$。其像是理想 $\mathcal I$，满足 $\mathcal L^{-1}\xrightarrow{\sim}\mathcal I$。因此 $\mathcal I$ 是某个有效除子 $D$ 的理想。于是 $\mathcal O_X(-D)=\mathcal I$，故 $\mathcal L^{-1}\xrightarrow{\sim}\mathcal O_X(-D)$。取逆得到 $\mathcal O_X(D)\simeq\mathcal L$。所以 $D\in|\mathcal L|$。因此有映射 $H^0(X,\mathcal L)_{\mathrm{reg}}\to|\mathcal L|$。

若截面乘以 $H^0(X,\mathcal O_X^*)$ 中的可逆元，则单射 $\mathcal L^{-1}\hookrightarrow\mathcal O_X$ 也乘以同一可逆元，因而像 $\mathcal I$ 相同，所以 $D$ 不变。反过来，若 $D$ 来自另一个截面，对应于另一个同构 $\mathcal L^{-1}\xrightarrow{\sim}\mathcal I$，则两个同构相差 $\mathcal L^{-1}$ 的自同构，它由乘以 $H^0(X,\mathcal O_X^*)$ 中的可逆元给出；因此两个截面相差乘以这个可逆元。于是
\[
 H^0(X,\mathcal L)_{\mathrm{reg}}/H^0(X,\mathcal O_X^*)\hookrightarrow|\mathcal L|.
\]
最后，给定 $D\in|\mathcal L|$，由定义存在同构 $\mathcal O_X(D)\simeq\mathcal L$。由于 $\mathcal O_X(-D)$ 是 $D$ 的理想 $\mathcal I$，包含 $\mathcal I\hookrightarrow\mathcal O_X$ 给出单射 $\mathcal L^{-1}\hookrightarrow\mathcal O_X$。后者对应于 $H^0(X,\mathcal L)_{\mathrm{reg}}$ 的截面，它按第一段的方法给出 $D$。因此
\[
 H^0(X,\mathcal L)_{\mathrm{reg}}/H^0(X,\mathcal O_X^*)\xrightarrow{\sim}|\mathcal L|.
\]
\hfill$\square$
\end{origstatement}

\begin{origstatement}{解答}{9.3.5}\label{appA:ans935}
设 $x\in D+E$。若 $x\notin D\cap E$，则 $D+E$ 在 $x$ 处是相对有效除子，因为 $D+E$ 在 $x$ 的某个邻域上等于 $D$ 或 $E$。所以假设 $x\in D\cap E$。引理 9.3.4 说明，$X$ 在 $x$ 处 $S$-平坦，且 $D,E$ 在 $x$ 处各由一个在通过 $x$ 的纤维 $X_s$ 上正则的元素切出。取这两个元素的乘积。
% Source PDF 315 = printed 303.
显然，它在 $x$ 处切出 $D+E$，且也在 $X_s$ 上正则。因此，再由引理 9.3.4，$D+E$ 在 $x$ 处是相对有效除子。\hfill$\square$
\end{origstatement}

\begin{origstatement}{解答}{9.3.8}\label{appA:ans938}
考虑 $X_T/T$ 上的相对有效除子 $D$。每个纤维 $D_t$ 的维数为 $0$，所以其 Hilbert 多项式 $\chi(\mathcal O_{D_t}(n))$ 为常数，值为 $\dim H^0(\mathcal O_{D_t})$，恰好是 $D_t$ 的次数。

各个断言在 $S$ 上局部，所以可假设 $X/S$ 射影。由定理 9.3.7，$\mathbf{Div}_{X/S}$ 由开子概形 $\mathbf{Div}_{X/S}\subseteq\mathbf{Hilb}_{X/S}$ 表示。而 $\mathbf{Hilb}_{X/S}$ 是有限型开闭子概形 $\mathbf{Hilb}_{X/S}^{\varphi}$ 的不交并，它们参数化 Hilbert 多项式为 $\varphi$ 的子概形。置
\[
 \mathbf{Div}_{X/S}^m\defeq\mathbf{Div}_{X/S}\cap\mathbf{Hilb}_{X/S}^m.
\]
那么 $\mathbf{Div}_{X/S}^m$ 具有所需的全部性质。

一般地，只要 $X/S$ 分离，$X$ 就表示 $\Hilb_{X/S}^1$，对角子概形 $\Delta\subset X\times X$ 是万有子概形。事实上，投影 $\Delta\to X$ 是同构，所以 $\Delta\in\Hilb_{X/S}^1(X)$。给定任意 $S$-映射 $g\colon T\to X$，注意 $(1\times g)^{-1}\Delta=\Gamma_g$，其中 $\Gamma_g\subset X\times T$ 是 $g$ 的图子概形，因为 $(1\times g)^{-1}\Delta$ 和 $\Gamma_g$ 的 $T'$-点恰好都是点对 $(gp,p)$，其中 $p\colon T'\to T$。因此 $\Gamma_g\in\Hilb_{X/S}^1(T)$。

反过来，设 $\Gamma\in\Hilb_{X/S}^1(T)$。那么 $\Gamma$ 是 $X\times T$ 的闭子概形。投影 $\pi\colon\Gamma\to T$ 固有，且纤维有限；因此由 Chevalley 定理 [EGAIII1, 4.4.2]，它是有限态射。所以 $\Gamma=\Spec(\pi_*\mathcal O_\Gamma)$。此外，$\pi_*\mathcal O_\Gamma$ 在 $\mathcal O_T$ 上平坦且有限生成，因而局部自由；其形成与取纤维可交换，所以秩为 $1$。因此 $\mathcal O_T\xrightarrow{\sim}\pi_*\mathcal O_\Gamma$，故 $\pi$ 是同构。于是 $\Gamma$ 是某个映射 $g\colon T\to X$ 的图。所以由上文，$(1\times g)^{-1}\Delta=\Gamma$；而 $g$ 是唯一具有此性质的映射，因为映射由其图确定。因此 $X$ 表示 $\Hilb_{X/S}^1$，且 $\Delta\subset X\times X$ 是万有子概形。

在这里的情形下，由定理 9.3.7，$\mathbf{Div}_{X/S}^1$ 因而可由开子概形 $U\subset X$ 表示。事实上，该定理的证明表明，$U$ 由满足纤维 $\Delta_x$ 为 $X_x$ 上除子的点 $x\in X$ 组成。对任意 $x\in X$，$\Delta_x$ 是 $k_x$-有理点，因此 $\Delta_x$ 是除子，当且仅当 $X_x$ 在 $\Delta_x$ 处正则。由于 $X/S$ 平坦，$X_x$ 在 $\Delta_x$ 处正则，当且仅当 $x\in X_0$。所以 $X_0=\mathbf{Div}_{X/S}^1$。

最后，置 $T\defeq X_0^m$，令 $\Gamma_i\subset X\times T$ 为第 $i$ 个投影的图子概形。由上述分析，$\Gamma_i\in\Div_{X/S}^1(T)$。置 $\Gamma\defeq\sum\Gamma_i$。由习题 9.3.5 及次数的可加性，$\Gamma\in\Div_{X/S}^m(T)$。显然 $\Gamma$ 表示所需的 $\mathbf{Div}_{X/S}^m$ 的 $T$-点。\hfill$\square$
\end{origstatement}

\begin{origstatement}{解答}{9.3.11}\label{appA:ans9311}
设 $s\in S$。令 $K$ 为 $k_s$ 的代数闭包，置 $A\defeq H^0(X_K,\mathcal O_{X_K})$。由于 $f$ 固有，$A$ 作为 $K$-向量空间有限维，所以 $A$ 是 Artin 环。由于 $X_K$ 连通，由 [EGAIII2, 7.8.6.1]，$A$ 不是两个非零环的积；所以 $A$ 是 Artin 局部环。由于 $X_K$ 约化，$A$ 约化，因此 $A$ 是域，是 $K$ 的有限扩张。又因 $K$ 代数闭，故 $A=K$。由于上同调与平坦基变换可交换，得到 $k_s\xrightarrow{\sim}H^0(X_s,\mathcal O_{X_s})$。

同构 $k_s\xrightarrow{\sim}H^0(X_s,\mathcal O_{X_s})$ 通过 $f_*(\mathcal O_X)\otimes k_s$ 分解为
\[
 k_s\longrightarrow f_*(\mathcal O_X)\otimes k_s\longrightarrow H^0(X_s,\mathcal O_{X_s}).
\]
% Source PDF 316 = printed 304.
所以第二个映射满射。对 $\mathcal F\defeq\mathcal O_X$、$\mathcal N\defeq k_s$，由小节 9.3.10 的蕴含 (iv)$\Rightarrow$(iii)，该映射是同构。因此第一个映射也是同构。

于是 $\mathcal O_S\to f_*\mathcal O_X$ 在 $s$ 处满射。事实上，以 $\mathcal G$ 表示其余核。由于张量积右正合，且 $k_s\to f_*(\mathcal O_X)\otimes k_s$ 是同构，有 $\mathcal G\otimes k_s=0$。因此由 Nakayama 引理，茎 $\mathcal G_s$ 消失，断言得证。

设 $\mathcal Q$ 是小节 9.3.10 中与 $\mathcal F\defeq\mathcal O_X$ 相伴的 $\mathcal O_S$-模。由小节 9.3.10 的蕴含 (iv)$\Rightarrow$(i)，$\mathcal Q$ 在 $s$ 处自由。并且，对 $\mathcal N\defeq k_s$，由 (9.3.10.1) 中的同构，$\operatorname{rank}\mathcal Q_s=1$。但对 $\mathcal N\defeq\mathcal O_S$，该同构成为 $\Hom(\mathcal Q,\mathcal O_X)\xrightarrow{\sim}f_*\mathcal O_X$。因此 $f_*\mathcal O_X$ 在 $s$ 处也自由且秩为 $1$。所以满射 $\mathcal O_S\to f_*\mathcal O_X$ 在 $s$ 处是同构。由于 $s$ 任意，$\mathcal O_S\xrightarrow{\sim}f_*\mathcal O_X$ 处处成立。

最后，设 $T$ 为任意 $S$-概形。那么 $f_T\colon X_T\to T$ 也固有且平坦，几何纤维约化且连通。因此由刚才所证，$\mathcal O_T\xrightarrow{\sim}f_{T*}\mathcal O_{X_T}$。\hfill$\square$
\end{origstatement}

\begin{origstatement}{解答}{9.3.14}\label{appA:ans9314}
由定理 9.3.13，$L$ 表示 $\mathrm{LinSys}_{\mathcal L/X/S}$。所以由 Yoneda 引理 [EGAG, (0,1.1.4), p.20]，存在 $W\in\mathrm{LinSys}_{\mathcal L/X/S}(L)$，具有所需的万有性质。$W$ 对应于恒等映射 $p\colon L\to L$。定理 9.3.13 的证明表明 $\mathcal O_{X_L}(W)=(\mathcal L|X_L)\otimes f_L^*\mathcal O_L(1)$。\hfill$\square$
\end{origstatement}

\begin{origstatement}{解答}{9.4.2}\label{appA:ans942}
结构层 $\mathcal O_X$ 定义截面 $\sigma\colon S\to\mathbf{Pic}_{X/S}$。其像是子概形；若 $\mathbf{Pic}_{X/S}$ 分离，则它为闭子概形，见 [EGAG, Cors.(5.1.4), p.275, 及 (5.2.4), p.278]。令 $N\subset T$ 是这个子概形在 $\mathcal L$ 定义的映射 $\lambda\colon T\to\mathbf{Pic}_{X/S}$ 下的拉回。那么第三个性质成立。

$\mathcal L_N$ 与 $\mathcal O_X$ 定义相同的映射 $N\to\mathbf{Pic}_{X/S}$。由于 $\mathcal O_S\xrightarrow{\sim}f_*\mathcal O_X$ 普遍成立，比较定理 9.2.5 蕴含，$N$ 上存在可逆层 $\mathcal N$，使第一个性质成立。

考虑第二个性质。此时 $\mathcal L_{T'}\simeq f_{T'}^*\mathcal N'$。所以 $\lambda t\colon T'\to\mathbf{Pic}_{X/S}$ 也由 $\mathcal O_{X_{T'}}$ 定义；因此 $\lambda t$ 通过 $\sigma\colon S\to\mathbf{Pic}_{X/S}$ 分解。故 $t\colon T'\to T$ 通过 $N$ 分解。由第一个性质，$\mathcal L_{T'}\simeq f_{T'}^*t^*\mathcal N$。由引理 9.2.7，得到 $\mathcal N'\simeq t^*\mathcal N$。所以第二个性质成立。

最后，假设对 $(N_1,\mathcal N_1)$，第一个性质也成立。取 $t$ 为 $N_1$ 到 $T$ 的包含，得到 $N_1\subset N$ 且 $\mathcal N_1\simeq\mathcal N|N$。若 $(N_1,\mathcal N_1)$ 也具有第二个性质，则类似地有 $N\subset N_1$。因此 $N=N_1$ 且 $\mathcal N_1\simeq\mathcal N$，正如所求。\hfill$\square$
\end{origstatement}

\begin{origstatement}{解答}{9.4.3}\label{appA:ans943}
由 Yoneda 引理 [EGAG, (0,1.1.4), p.20]，万有层 $\mathcal P$ 存在，当且仅当 $\mathbf{Pic}_{X/S}$ 表示 $\Pic_{X/S}$。置 $P\defeq\mathbf{Pic}_{X/S}$。

假设 $\mathcal P$ 存在。那么对 $P$ 上任意可逆层 $\mathcal N$，显然 $\mathcal P\otimes f_P^*\mathcal N$ 也是万有层。若 $\mathcal P'$ 也是万有层，则在定义中取 $h\defeq1_P$，得到对 $P$ 上某个可逆层 $\mathcal N$，有 $\mathcal P'\simeq\mathcal P\otimes f_P^*\mathcal N$。

假设 $\mathcal O_S\xrightarrow{\sim}f_*\mathcal O_X$ 普遍成立。若对 $P$ 上某些可逆层 $\mathcal N,\mathcal N'$，有 $\mathcal P\otimes f_P^*\mathcal N\simeq\mathcal P\otimes f_P^*\mathcal N'$，则由引理 9.2.7，$\mathcal N\simeq\mathcal N'$。

由定理 9.2.5 的第二部分，若 $f$ 还有截面，则 $\mathbf{Pic}_{X/S}$ 确实表示 $\Pic_{X/S}$，所以 $\mathcal P$ 存在。此外，习题 9.2.4 中的曲线 $X/\mathbb R$ 给出 $\mathcal P$ 不存在的例子，因为由定理 9.4.8，$\Pic_{(X/\mathbb R)\,(\text{ét})}$ 可表，而 $\Pic_{X/\mathbb R}$ 不可表，因为两个函子不同。\hfill$\square$
\end{origstatement}

% Source PDF 317 = printed 305.
\begin{origstatement}{解答}{9.4.4}\label{appA:ans944}
设 $\mathbf{Pic}_{X/S}$ 表示 $\Pic_{(X/S)\,(\text{ét})}$。对任意 $S'$-概形 $T$，有
\[
 \Pic_{(X_{S'}/S')\,(\text{ét})}(T)=\Pic_{(X/S)\,(\text{ét})}(T).
\]
这本质上由定义成立，因为 $S'$-概形的映射是平展覆盖，当且仅当它作为 $S$-概形的映射是平展覆盖。然而，由于结构映射 $T\to S'$ 固定，有
\[
 (\mathbf{Pic}_{X/S}\times_S S')(T)=\mathbf{Pic}_{X/S}(T).
\]
两个等式的右侧相等，所以左侧也相等。因此 $\mathbf{Pic}_{X/S}\times_S S'$ 表示 $\Pic_{(X_{S'}/S')\,(\text{ét})}$。当然，当 $\mathbf{Pic}_{X/S}$ 表示其余某个相对 Picard 函子时，也可作类似分析。

习题 9.2.4 中的曲线 $X\subset\mathbb P_{\mathbb R}^2$ 给出一个例子。事实上，因为函子 $\Pic_{X/\mathbb R}$ 与 $\Pic_{(X/\mathbb R)\,(\text{ét})}$ 不同，$\Pic_{X/\mathbb R}$ 不可表。但由主要定理 9.4.8，$\Pic_{(X/\mathbb R)\,(\text{ét})}$ 可表。最后，由于 $X_{\mathbb C}$ 有 $\mathbb C$-点，由比较定理 9.2.5，其所有相对 Picard 函子相等。\hfill$\square$
\end{origstatement}

\begin{origstatement}{解答}{9.4.5}\label{appA:ans945}
$X_k$ 上的 $\mathcal L$ 定义映射 $\Spec(k)\to\mathbf{Pic}_{X/S}$；把它的像对应于 $\mathcal L$。那么对包含 $k$ 的任意域 $k''$，拉回 $\mathcal L|X_{k''}$ 对应于 $\mathbf{Pic}_{X/S}$ 的同一个概形点。

考虑 $X_{k'}$ 上的 $\mathcal L'$。若 $\mathcal L,\mathcal L'$ 表示同一个类，则存在同时包含 $k,k'$ 的 $k''$，使 $\mathcal L|X_{k''}\simeq\mathcal L'|X_{k''}$；因此 $\mathcal L,\mathcal L'$ 都对应于 $\mathbf{Pic}_{X/S}$ 的同一个概形点。反过来，若 $\mathcal L,\mathcal L'$ 对应于同一点，取同时包含 $k,k'$ 的任意代数闭域 $k''$。那么 $\mathcal L|X_{k''}$ 与 $\mathcal L'|X_{k''}$ 定义相同的映射 $\Spec(k'')\to\mathbf{Pic}_{X/S}$，从而由习题 9.2.3 或 9.2.6，$\mathcal L|X_{k''}\simeq\mathcal L'|X_{k''}$。

最后，给定 $\mathbf{Pic}_{X/S}$ 的任意概形点，令 $k$ 是其剩余域的代数闭包。由习题 9.2.3 或 9.2.6，$\Spec(k)\to\mathbf{Pic}_{X/S}$ 由 $X_k$ 上某个 $\mathcal L$ 定义。所以给定点对应于 $\mathcal L$。因此，$X/S$ 的纤维上可逆层的类与 $\mathbf{Pic}_{X/S}$ 的概形点一一对应。\hfill$\square$
\end{origstatement}

\begin{origstatement}{解答}{9.4.7}\label{appA:ans947}
$S$-映射 $h\colon T\to\mathbf{Div}_{X/S}$ 对应于 $X_T$ 上的相对有效除子 $D$。因此复合 $\mathbf A_{X/S}h\colon T\to P$ 对应于可逆层 $\mathcal O_{X_T}(D)$。所以对 $T$ 上某个可逆层 $\mathcal N$，有 $\mathcal O_{X_T}(D)\simeq(1\times\mathbf A_{X/S}h)^*\mathcal P\otimes f_T^*\mathcal N$。故若通过 $\mathbf A_{X/S}h$ 把 $T$ 看作 $P$-概形，$D$ 就定义 $\mathrm{LinSys}_{\mathcal P/X\times P/P}$ 的 $T$-点 $\eta$。显然，对应 $h\mapsto\eta$ 对 $T$ 具有函子性。因此，若通过 $\mathbf A_{X/S}$ 把 $\mathbf{Div}_{X/S}$ 看作 $P$-概形，就有从其点函子到 $\mathrm{LinSys}_{\mathcal P/X\times P/P}$ 的自然映射 $\Lambda$。

此外，$\Lambda$ 是同构。事实上，设 $T$ 是 $P$-概形。$\mathrm{LinSys}_{\mathcal P/X\times P/P}$ 的 $T$-点 $\eta$ 由 $X_T$ 上的相对有效除子 $D$ 给出，满足对 $T$ 上某个可逆层 $\mathcal N$，$\mathcal O_{X_T}(D)\simeq\mathcal P_T\otimes f_T^*\mathcal N$。那么 $\mathcal O_{X_T}(D)$ 与 $\mathcal P_T$ 定义相同的 $S$-映射 $T\to P$。但 $\mathcal P_T$ 定义结构映射，而 $\mathcal O_{X_T}(D)$ 定义复合 $\mathbf A_{X/S}h$，其中 $h\colon T\to\mathbf{Div}_{X/S}$ 由 $D$ 定义。因此 $\eta=\Lambda(h)$，且 $h$ 由 $\eta$ 确定，故 $\Lambda$ 是同构。

换言之，$\mathbf{Div}_{X/S}$ 表示 $\mathrm{LinSys}_{\mathcal P/X\times P/P}$。但由定理 9.3.13，$\mathbb P(\mathcal Q)$ 也表示 $\mathrm{LinSys}_{\mathcal P/X\times P/P}$。因此作为 $P$-概形，$\mathbb P(\mathcal Q)=\mathbf{Div}_{X/S}$。\hfill$\square$
\end{origstatement}

% Source PDF 318 = printed 306.
\begin{origstatement}{解答}{9.4.10}\label{appA:ans9410}
先假设 $F\to G$ 是满射。给定平展层的映射 $\varphi\colon F\to H$，使两个映射 $F\times_G F\to H$ 相等，必须证明存在唯一映射 $G\to H$，使 $F\to G\to H$ 等于 $\varphi$。

设 $\eta\in G(T)$。由假设，存在平展覆盖 $T'\to T$ 及元素 $\zeta'\in F(T')$，使 $\zeta'$ 与 $\eta$ 在 $G(T')$ 中的像相同。置 $T''\defeq T'\times_T T'$。那么 $\zeta'$ 在 $F(T'')$ 中的两个像定义 $(F\times_G F)(T'')$ 的元素 $\zeta''$。由于两个映射 $F\times_G F\to H$ 相等，$\zeta''$ 在 $H(T'')$ 中的两个像相等。但这两个像恰好是 $\varphi(\zeta')\in H(T')$ 的两个像。由于 $H$ 是层，$\varphi(\zeta')$ 是唯一元素 $\theta\in H(T)$ 的像。

注意，$\theta\in H(T)$ 不依赖于 $T'$ 和 $\zeta'\in F(T')$ 的选择。事实上，设 $\zeta'_1\in F(T'_1)$ 是另一种选择。按上面的论证，$\varphi(\zeta'_1)\in H(T'_1)$ 和 $\varphi(\zeta')\in H(T')$ 在 $H(T'_1\times_T T')$ 中有相同的像。所以 $\zeta'_1$ 也给出 $\theta$。

以 $\eta\mapsto\theta$ 定义映射 $G(T)\to H(T)$。显然，它对 $T$ 具有函子性，因而有层映射 $G\to H$。显然 $F\to G\to H$ 等于 $\varphi\colon F\to H$。最后，$G\to H$ 是唯一这样的映射，因为 $\eta$ 在 $H(T)$ 中的像由它在 $G(T')$ 中的像确定，而后者必须映到 $\varphi(\zeta')\in H(T')$。因此 $G$ 是 $F\times_G F\rightrightarrows F$ 的余等化子。

反过来，假设 $G$ 是 $F\times_G F\rightrightarrows F$ 的余等化子。取平展子层 $H\subset G$，它相伴于下述预层：其 $T$-点是 $F(T)$ 的元素在 $G(T)$ 中的像。于是映射 $F\to G$ 通过 $H$ 分解，所以两个映射 $F\times_G F\to H$ 相等。由于 $G$ 是余等化子，存在映射 $G\to H$，使 $F\to G\to H$ 等于 $F\to H$。因此 $F\to G\to H\hookrightarrow G$ 等于 $F\to G$。由唯一性，$G\to H\hookrightarrow G$ 等于 $1_G$。故 $H=G$。所以 $F\to G$ 是满射。\hfill$\square$
\end{origstatement}

\begin{origstatement}{解答}{9.4.11}\label{appA:ans9411}
由定理 9.4.8，$Z$ 的每个连通分支 $Z'$ 包含在 $\mathbf{Pic}_{X/S}$ 的拟射影开子概形的递增并中。由于 $Z'$ 拟紧，$Z'$ 包含在其中某一个子概形中，所以 $Z'$ 拟射影。但 $Z$ 只有有限多个分支 $Z'$。因此 $Z$ 拟射影。\hfill$\square$
\end{origstatement}

\begin{origstatement}{解答}{9.4.12}\label{appA:ans9412}
置 $P\defeq\mathbf{Pic}_{X/S}$，由定理 9.4.8 它存在。若 $\mathcal P$ 存在，则由习题 9.4.7，$\mathbf A_{X/S}$ 是丛 $\mathbb P(\mathcal Q)$ 的结构映射，其中 $\mathcal Q$ 是小节 9.3.10 中与 $\mathcal P$ 相伴的 $\mathbf{Pic}_{X/S}$ 上的凝聚层。特别地，$\mathbf A_{X/S}$ 在 $S$ 上 Zariski 局部地射影。

一般地，由习题 9.4.4，$P$ 的形成与扩张 $S$ 可交换。类似地，$\mathbf A_{X/S}$ 的形成也是如此。而由 [EGAIV2, 2.7.1(vii)]，映射若在 fppf 基扩张后固有，就本身固有。

然而，$f\colon X\to S$ 是 fppf 态射。并且 $f_X\colon X\times X\to X$ 有截面，即对角态射。因此用 $f$ 作基扩张。由习题 9.3.11 和 9.4.3，存在万有层 $\mathcal P$。于是由第一种情形，$\mathbf A_{X/S}$ 固有。\hfill$\square$
\end{origstatement}

% Source PDF318, printed306, final answer begins here and continues onto PDF319.
\begin{origstatement}{解答}{9.4.13}\label{appA:ans9413}
采用解答 9.4.12 的思想和记号。由习题 9.3.8，$X_0$ 表示 $\Div_{X/S}^1$。因此 Abel 映射 $\mathbf A_{X/S}$ 诱导自然映射 $A\colon X_0\to P$，且 $A$ 的形成与扩张 $S$ 可交换。由 [EGAIV2, 2.7.1(xii)]，一个映射若在 fppf 基扩张后为闭浸入，就本身为闭浸入。因此可假设 $\mathbb P(\mathcal Q)=\mathbf{Div}_{X/S}$。

函数 $\lambda\mapsto\deg\mathcal P_\lambda$ 局部常值。设 $W\subset P$ 为该函数取值 $1$ 的开闭子集。由上文，显然 $\mathbb P(\mathcal Q_W)=\mathbf{Div}_{X/S}^1$。
% Source PDF 319 = printed 307.
因此 $X_0=\mathbb P(\mathcal Q_W)$，而 $A\colon X_0\to P$ 等于 $\mathbb P(\mathcal Q_W)$ 的结构映射。所以还需证明这个结构映射是闭浸入。

固定 $\lambda\in W$。有 $\dim_{k_\lambda}(\mathcal Q\otimes k_\lambda)=\dim_{k_\lambda}H^0(X_\lambda,\mathcal P_\lambda)$。假设 $\mathcal P_\lambda$ 有两个独立的整体截面。每个截面都定义次数为 $1$ 的有效除子，即一个 $k_\lambda$-有理点 $x_i$。由于两个截面互不为倍数，$x_i$ 互不相同。因此这些截面生成 $\mathcal P_\lambda$，从而由 [EGAII, 4.2.3] 或 [Har77, Thm.II, 7.1, p.150]，定义映射 $h\colon X_\lambda\to\mathbb P_{k_\lambda}^1$。$h$ 双有理，因为每个 $x_i$ 都是 $\mathbb P_{k_\lambda}^1$ 的一个 $k_\lambda$-有理点的概形论逆像。因此 $h$ 是同构。但由假设，$X_\lambda$ 的算术亏格至少为 $1$，矛盾。所以 $\dim_{k_\lambda}(\mathcal Q\otimes k_\lambda)\leq1$。

由 Nakayama 引理，$\mathcal Q$ 在 $\lambda$ 的某个邻域 $V\subset W$ 上可由一个元素生成。因此有满射 $\mathcal O_V\twoheadrightarrow\mathcal Q_V$，定义闭浸入 $\mathbb P(\mathcal Q_V)\hookrightarrow\mathbb P(\mathcal O_V)$。但结构映射 $\mathbb P(\mathcal O_V)\to V$ 是同构，因此 $\mathbb P(\mathcal Q_V)\to V$ 是闭浸入。由于 $\lambda\in W$ 任意，$\mathbb P(\mathcal Q_W)\to W$ 的确是闭浸入。\hfill$\square$
\end{origstatement}

\begin{origstatement}{解答}{9.4.15}\label{appA:ans9415}
表示 $\Pic_{X/S}$ 与例 9.4.14 中表示 $\Pic_{X'/S'}$ 类似，但更简单。事实上，在 $X\times_S\mathbb Z_S$ 上，把 $\mathcal O_X(n)$ 放在第 $n$ 份 $X$ 上，构造可逆层 $\mathcal P$。那么只需证明：给定任意 $S$-概形 $T$ 和 $X_T$ 上任意可逆层 $\mathcal L$，存在唯一的 $S$-映射 $q\colon T\to\mathbb Z_S$ 以及 $T$ 上某个可逆层 $\mathcal N$，使 $(1\times q)^*\mathcal P=\mathcal L\otimes f_T^*\mathcal N$。

显然可假设 $T$ 连通。由 [EGAIII2, 7.9.11]，函数 $s\mapsto\chi(X_t,\mathcal L_t)$ 在 $T$ 上为常数。$X_t$ 是剩余域 $k_t$ 上维数至少为 $1$ 的射影空间，因此由 [Har77, Prp.6.4, p.132, 及 Cors.6.16 和 6.17, p.145]，对某个 $n$，$\mathcal L_t\simeq\mathcal O_{X_t}(n)$。所以 $n$ 不依赖于 $t$。

置 $\mathcal M\defeq\mathcal L^{-1}(n)$。那么对所有 $t\in T$，$\mathcal M_t\simeq\mathcal O_{X_t}$。由 Serre 的显式计算 [EGAIII1, 2.1.12]，$H^1(X_t,\mathcal M_t)=0$ 且 $H^0(X_t,\mathcal M_t)=k_t$。因此由小节 9.3.10 的理论，$f_{T*}\mathcal M$ 可逆，且其形成与改变基 $T$ 可交换。

置 $\mathcal N\defeq f_{T*}\mathcal M$，考虑自然映射 $u\colon f_T^*\mathcal N\to\mathcal M$。由于 $\mathcal N$ 的形成与改变 $T$ 可交换，$u$ 的形成也如此。但 $u$ 在每个 $t\in T$ 上的纤维上是同构，所以 $u\otimes k_t$ 是同构。由 Nakayama 引理，$u$ 满射。而 $u$ 的定义域和余域都可逆，因此 $u$ 是同构。故 $\mathcal L\otimes f_T^*\mathcal N=\mathcal O_{X_T}(n)$。

令 $q\colon T\to\mathbb Z_S$ 为结构映射 $T\to S$ 与第 $n$ 个包含 $S\hookrightarrow\mathbb Z_S$ 的复合。显然 $(1\times q)^*\mathcal P=\mathcal O_{X_T}(n)$，且 $q$ 是唯一这样的 $S$-映射。因此 $\mathbb Z_S$ 表示 $\Pic_{X/S}$，而 $\mathcal P$ 是万有层。\hfill$\square$
\end{origstatement}

\begin{origstatement}{解答}{9.4.16}\label{appA:ans9416}
首先，由定理 9.4.8，$\mathbf{Pic}_{X/\mathbb R}$ 存在。$X_{\mathbb C}\simeq\mathbb P_{\mathbb C}^1$，因此由习题 9.4.4 和 9.4.15，$\mathbf{Pic}_{X/\mathbb R}\times_{\mathbb R}\mathbb C\simeq\mathbb Z_{\mathbb C}$。在 $\mathbb Z_{\mathbb C\otimes_{\mathbb R}\mathbb C}$ 上诱导的自同构是恒等映射；事实上，该概形的点对应于 $\mathbb P_{\mathbb C}^1$ 上的可逆层，每个这样的层都同构于它在 $\mathbb P_{\mathbb C}^1$ 的任意 $\mathbb R$-自同构下的拉回。因此由下降理论，$\mathbf{Pic}_{X/\mathbb R}=\mathbb Z_{\mathbb R}$。

上述论证还给出 $\Pic_{(X/\mathbb R)\,(\text{ét})}$ 可表的第二个证明。事实上，置 $P\defeq\Pic_{(X/\mathbb R)\,(\text{ét})}$。按上面的论证，映射对
\[
 (P\otimes_{\mathbb R}\mathbb C)\otimes_{\mathbb C}(\mathbb C\otimes_{\mathbb R}\mathbb C)
 \rightrightarrows P\otimes_{\mathbb R}\mathbb C
\]
% Source PDF 320 = printed 308.
由映射对 $\mathbb Z_{\mathbb C\otimes_{\mathbb R}\mathbb C}\rightrightarrows\mathbb Z_{\mathbb C}$ 表示，其余等化子是 $\mathbb Z_{\mathbb R}$。另一方面，由习题 9.4.10，在平展层范畴中，其余等化子是 $P$。

顺便注意，$\mathbf{Pic}_{X/\mathbb R}=\mathbf{Pic}_{\mathbb P_{\mathbb R}^1/\mathbb R}$。但由习题 9.2.4，$\Pic_{X/\mathbb R}$ 不可表，而由习题 9.4.15，$\Pic_{\mathbb P_{\mathbb R}^1/\mathbb R}$ 可表。\hfill$\square$
\end{origstatement}

\begin{origstatement}{解答}{9.5.7}\label{appA:ans957}
习题 9.4.11 蕴含 $Z$ 拟射影。因此若 $Z$ 固有，便射影。由 [EGAIV2, 2.7.1]，一个 $S$-概形若在如 $f\colon X\to S$ 这样的 fppf 基变换后固有，就本身固有。而 $f_X\colon X\times X\to X$ 有截面，即对角态射。因此可假设 $f$ 有截面。

使用固有性的赋值判别法 [Har77, Thm.4.7, p.101]，只需验证：给定形如 $T=\Spec(A)$ 的 $S$-概形 $T$，其中 $A$ 是赋值环，分式域为 $K$，每个 $S$-映射 $u\colon\Spec(K)\to Z$ 都扩张为 $S$-映射 $T\to Z$。不必验证扩张若存在是否唯一；事实上，由于 $Z$ 拟射影，因而分离，分离性的赋值判别法 [Har77, Thm.4.3, p.97] 已保证该唯一性。

由于 $f$ 有截面，由定理 9.2.5，$u$ 来自 $X_K$ 上的可逆层 $\mathcal L$。需要把 $\mathcal L$ 扩张到 $X_T$ 上。这样的扩张定义映射 $t\colon T\to\mathbf{Pic}_{X/S}$，扩张 $u$；而由于 $Z$ 闭且 $T$ 整，$t$ 通过 $Z$ 分解。

显然，只需对任意 $n\gg0$ 扩张 $\mathcal L(n)$。因此必要时替换 $\mathcal L$，可假设 $\mathcal L$ 有非零截面。由于 $X_K$ 整，该截面正则。因此 $X_K$ 上有除子 $D$，使 $\mathcal O(D)=\mathcal L$。

设 $D'\subset X_T$ 为 $D$ 的闭包。由于 $X/S$ 光滑且 $T$ 正则，由 [EGAIV2, 6.5.2]，$X_T$ 正则，从而由 [EGAIV2, 21.11.1]，它是析因的。因此 $D'$ 是除子，$\mathcal O(D')$ 扩张 $\mathcal L$。\hfill$\square$
\end{origstatement}

\begin{origstatement}{解答}{9.5.16}\label{appA:ans9516}
由 Serre 定理 [Har77, Thm.5.2, p.228]，当 $i>0$ 且 $n\gg0$ 时，$H^i(\Omega_X^2(n))=0$。所以 $\varphi(n)=\chi(\Omega_X^2(n))$。于是
\[
 q=H^1(\Omega_X^2)-H^2(\Omega_X^2)+1.
\]
Serre 对偶 [Har77, Cor.7.13, p.247] 给出，对一切 $i$，$\dim H^i(\Omega_X^2)=\dim H^{2-i}(\mathcal O_X)$。又由于 $X$ 射影且几何整，$\dim H^0(\mathcal O_X)=1$。所以
\[
 q=\dim H^1(\mathcal O_X).
\]
因此由推论 9.5.14，$\dim\mathbf{Pic}_{X/S}\leq q$，且在特征 $0$ 时取等号。\hfill$\square$
\end{origstatement}

\begin{origstatement}{解答}{9.5.17}\label{appA:ans9517}
置 $P\defeq\mathbf{Pic}_{X/S}$，由定理 9.4.8 它存在。由习题 9.3.11 和 9.4.3，$X\times P$ 上存在万有层 $\mathcal P$。

假设 $q=0$。那么由推论 9.5.13，$P$ 处处光滑且维数为 $0$。设 $D$ 是 $X_T/T$ 上的相对有效除子，其中 $T$ 是连通 $S$-概形。于是 $\mathcal O_{X_T}(D)$ 定义映射 $\tau\colon T\to P$，且对 $T$ 上某个可逆层 $\mathcal N$，有
\[
 \mathcal O_{X_T}(D)\simeq(1\times\tau)^*\mathcal P\otimes f_T^*\mathcal N.
\]
$T$ 连通而 $P$ 离散且约化，所以 $\tau$ 为常值。置 $\lambda\defeq\tau T$，把 $\mathcal P_\lambda$ 看作 $X$ 上的可逆层 $\mathcal L$。那么 $\mathcal L_T=(1\times\tau)^*\mathcal P$，因此 $\mathcal O_{X_T}(D)\simeq\mathcal L_T\otimes f_T^*\mathcal N$，正如所求。

考虑逆命题。再由习题 9.4.3，$P$ 上有凝聚层 $\mathcal Q$，使 $\mathbb P(\mathcal Q)=\mathbf{Div}_{X/S}$。并且，由小节 9.3.10，$\mathcal Q$ 在任何表示 $X$ 上满足 $H^1(\mathcal L)=0$ 的可逆层 $\mathcal L$ 的闭点 $\lambda$ 处都非零且局部自由，
% Source PDF 321 = printed 309.
例如，可取 $\mathcal L\defeq\mathcal O_X(n)$，其中 $n\gg0$。

设 $U\subset P$ 是 $\lambda$ 的连通开邻域，使 $\mathcal Q$ 在其上自由。设 $T\subset\mathbb P(\mathcal Q)$ 是 $U$ 的原像，$D$ 是 $X_T/T$ 上的万有相对有效除子。那么自然映射 $A\colon T\to U$ 光滑且纤维不可约，所以 $T$ 连通。此外，$A$ 是 $\mathcal O_{X_T}(D)$ 定义的映射。

假设对 $X$ 上某个可逆层 $\mathcal M$ 及 $T$ 上某个可逆层 $\mathcal N$，有 $\mathcal O_{X_T}(D)\simeq\mathcal M_T\otimes f_T^*\mathcal N$。那么 $A\colon T\to U$ 也由 $\mathcal M_T$ 定义。设 $\mu\in P$ 表示 $\mathcal M$，则 $A$ 通过闭点 $\mu$ 的包含分解。于是 $\mu=\lambda$；且由于 $A$ 光滑满射，其像，即开集 $U$，就是约化闭点 $\lambda$。$P$ 有一个把 $0$ 送到 $\lambda$ 的自同构，即与 $\lambda$“相乘”。所以 $P$ 在 $0$ 处光滑且维数为 $0$。由推论 9.5.13，$q=0$。

在特征 $0$ 时，事先由推论 9.5.14 已知 $P$ 光滑。$A\colon T\to U$ 光滑，因此 $T$ 也光滑。但前面的论证表明，若条件对这个 $T$ 成立，则 $q=0$，正如所求。\hfill$\square$
\end{origstatement}

\begin{origstatement}{解答}{9.5.23}\label{appA:ans9523}
由假设，对 $s\in S$，有 $\dim X_s=1$，因此 $H^2(\mathcal O_{X_s})=0$。由命题 9.5.19，$\mathbf{Pic}_{X_s/k_s}^0$ 光滑，因而由命题 9.5.13，其维数为 $p_a$。因此由命题 9.5.20，$\mathbf{Pic}_{X_s/k_s}^0$ 构成有限型族，其全空间是 $\mathbf{Pic}_{X/S}$ 的开子概形 $\mathbf{Pic}_{X/S}^0$。再由命题 9.5.19，$\mathbf{Pic}_{X/S}$ 在 $S$ 上光滑。

因此由习题 9.4.11，$\mathbf{Pic}_{X/S}^0$ 拟射影。

若 $X/S$ 光滑，则由习题 9.5.7，$\mathbf{Pic}_{X/S}^0$ 在 $S$ 上射影。也可再用定理 9.5.4 和命题 9.5.20，得出 $\mathbf{Pic}_{X/S}^0$ 固有，因而由于拟射影而射影。

反过来，假设 $\mathbf{Pic}_{X/S}^0$ 固有，证明 $X/S$ 光滑。由于 $X/S$ 平坦，只需证明每个 $X_s$ 光滑。因此可用 $k_s$ 的代数闭包的谱代替 $S$。若 $p_a=0$，则 $X$ 光滑，事实上由 [Har77, Ex.1.8(b), p.298]，$X=\mathbb P^1$。

假设 $p_a>0$。设 $X_0$ 是 $X$ 光滑的开子概形。由习题 9.4.13，有闭浸入 $A\colon X_0\hookrightarrow\mathbf{Pic}_{X/S}$。其像由表示次数为 $1$ 的可逆层的点 $\lambda$ 组成。固定一个有理点 $\lambda$，定义 $\mathbf{Pic}_{X/S}$ 的自同构 $\beta$ 为 $\beta(\kappa)\defeq\kappa\lambda^{-1}$。那么 $\beta A$ 是 $X_0$ 到 $\mathbf{Pic}_{X/S}^0$ 的闭浸入。

由假设，$\mathbf{Pic}_{X/S}^0$ 固有，所以 $X_0$ 固有。由于 $X$ 分离，$X_0\hookrightarrow X$ 固有，因而 $X_0$ 在 $X$ 中闭。但由于 $X$ 整且底域代数闭，$X_0$ 在 $X$ 中稠密。因此 $X_0=X$，也就是说 $X$ 光滑。\hfill$\square$
\end{origstatement}

\begin{origstatement}{解答}{9.6.4}\label{appA:ans964}
如前所述，由引理 9.6.6，存在 $m$，使每个 $\mathcal N(m)$ 都由整体截面生成。因此，给定 $X$ 的任意伴随点，都有一个截面不在该点消失；由于伴随点只有有限个，取相应截面的一般线性组合 $\sigma$，则 $\sigma$ 在任何伴随点都不消失。所以 $\sigma$ 正则，从而定义有效除子 $D$，使
\[
 \mathcal O_X(-D)=\mathcal N^{-1}(-m).
\]
% Source PDF 322 = printed 310.
显然 $\mathcal N^{-1}$ 也数值等价于 $\mathcal O_X$，所以由引理 9.6.6，$\chi(\mathcal N^{-1}(n))=\chi(\mathcal O_X(n))$。因此序列 $0\to\mathcal O_X(-D)\to\mathcal O_X\to\mathcal O_D\to0$ 给出
\[
 \chi(\mathcal O_D(n))=\psi(n),\qquad
 \psi(n)\defeq\chi(\mathcal O_X(n))-\chi(\mathcal O_X(n-m)).
\]
设 $T\subset\mathbf{Div}_{X/k}$ 是参数化 Hilbert 多项式为 $\psi(n)$ 的有效除子的开闭子概形。那么 $T$ 是有限型 $k$-概形。设 $\mathcal M'$ 是 $X_T$ 上与万有除子相伴的可逆层，置 $\mathcal M\defeq\mathcal M'(-n)$。那么存在有理点 $t\in T$，使 $\mathcal N=\mathcal M_t$。因此，数值等价于 $\mathcal O_X$ 的 $\mathcal N$ 构成有界族。\hfill$\square$
\end{origstatement}

\begin{origstatement}{解答}{9.6.7}\label{appA:ans967}
假设 $a<a_r$。假设 $\mathcal L(-1)$ 有非零截面，它定义有效除子 $D$，可能为 $0$。于是
\[
 0\leq\int h^{r-1}[D]=\int h^{r-1}\ell-\int h^r=a-a_r<0,
\]
矛盾。因此 $H^0(\mathcal L(-1))=0$。

设 $H$ 是 $X$ 的超平面截面。则有正合序列
\[
 0\longrightarrow\mathcal L(n-1)\longrightarrow\mathcal L(n)\longrightarrow\mathcal L_H(n)\longrightarrow0.
\]
它给出下述界：
\begin{equation}\tag{A.9.6.7.1}\label{appA:eqA9671}
 \dim H^0(\mathcal L(n))-\dim H^0(\mathcal L(n-1))\leq\dim H^0(\mathcal L_H(n)).
\end{equation}
由于 $\binom{n+i}{i}-\binom{n-1+i}{i}=\binom{n+i-1}{i-1}$，该序列还给出公式
\[
 \chi(\mathcal L_H(n))=\sum_{0\leq i\leq r-1}a_{i+1}\binom{n+i}{i}.
\]
假设 $r=1$。那么 $\dim H^0(\mathcal L_H(n))=\chi(\mathcal L_H(n))=a_1$。因此，由式 \eqref{appA:eqA9671} 对 $n$ 归纳，得到 $\dim H^0(\mathcal L(n))\leq a_1(n+1)$，正如所求。

此外，$\mathcal L_H$ 是 $0$-正则的。置 $m\defeq\dim H^1(\mathcal L(-1))$。由 [Mum66, p.102] 底部 Mumford 的结论，$\mathcal L$ 是 $m$-正则的。但
\[
 m=\dim H^0(\mathcal L(-1))-\chi(\mathcal L(-1))=0-a_1(-1+1)-a_0=-a_0.
\]
因此可取 $\Phi_1(u_0)\defeq-u_0$，其中 $u_0$ 为不定元。

假设 $r\geq2$。由 Bertini 定理 [Sei50, Thm.12, p.374] 或 [Jou79, Cor.6.7, p.80]，可取 $H$ 不可约。置 $h_1\defeq c_1\mathcal O_H(1)$ 及 $\ell_1\defeq c_1\mathcal L_H$。那么
$\int\ell_1h_1^{r-2}=\int\ell h^{r-2}[H]=a<a_r$。因此对 $r$ 归纳，可假设
\[
 \dim H^0(\mathcal L_H(n))\leq a_r\binom{n+r-1}{r-1}.
\]
于是由式 \eqref{appA:eqA9671} 对 $n$ 归纳，得到所需的界。

此外，可假设 $\mathcal L_H$ 是 $m_1$-正则的，其中 $m_1\defeq\Phi_{r-1}(a_1,\ldots,a_{r-1})$。置 $m\defeq m_1+\dim H^1(\mathcal L(m_1-1))$。由 Mumford 的同一工作，$\mathcal L$ 是 $m$-正则的。但
\[
\begin{aligned}
 m&=m_1+\dim H^0(\mathcal L(m_1-1))-\chi(\mathcal L(m_1-1))\\
 &\leq m_1+a_r\binom{m_1-1+r}{r}-\sum_{0\leq i\leq r}a_i\binom{m_1-1+i}{i}.
\end{aligned}
\]
后一个表达式是 $a_0,\ldots,a_{r-1}$ 和 $m_1$ 的多项式，所以它是只含 $a_0,\ldots,a_{r-1}$ 的多项式 $\Phi_r$，正如所求。

一般地，考虑 $\mathcal N\defeq\mathcal L(-a)$。那么
\[
 \chi(\mathcal N(n))=\sum_{0\leq i\leq r}b_i\binom{n+i}{i},\qquad
 b_i\defeq\sum_{j=0}^{r-i}a_{i+j}(-1)^j\binom{a-i-j}{j}.
\]
置 $\nu\defeq c_1\mathcal N$ 及 $b\defeq\int\nu h^{r-1}$。则
\[
 b=\int\ell h^{r-1}-a\int h^r=a-aa_r\leq0<a_r.
\]
% Source PDF 323 = printed 311.
因此 $\mathcal N$ 是 $m$-正则的，其中 $m\defeq\Phi_r(b_0,\ldots,b_{r-1})$。但 $b_i$ 是 $a_0,\ldots,a_r$ 及 $a$ 的多项式。因此存在只依赖于 $r$ 的多项式 $\Psi_r$，使 $m\defeq\Psi_r(a_0,\ldots,a_r;a)$，正如所求。\hfill$\square$
\end{origstatement}

\begin{origstatement}{解答}{9.6.10}\label{appA:ans9610}
令 $k'$ 为 $k$ 的代数闭包。若 $H\otimes k'\subset G^\tau\otimes k'$，则 $H\subset G^\tau$。但由引理 9.6.10，$G^\tau\otimes k'=(G\otimes k')^\tau$。因此可假设 $k=k'$。

此时 $H\subset\bigcup_{h\in H(k)}hG^0$。但 $G^0$ 开，所以 $hG^0$ 也开，而 $H$ 拟紧。因此 $H$ 包含在有限多个 $hG^0$ 中。所以 $G^0(k)$ 在 $H(k)G^0(k)$ 中的指数有限，记为 $n$。那么对每个 $h\in H(k)$，有 $h^n\in G^0(k)$。所以 $\varphi_n(H)\subset G^0$。于是 $H\subset G^\tau$。\hfill$\square$
\end{origstatement}

\begin{origstatement}{解答}{9.6.11}\label{appA:ans9611}
给定 $n$，显然 $\mathcal L^{\otimes n}$ 对应于 $\varphi_n\lambda$。由命题 9.5.10，$\mathcal L^{\otimes n}$ 代数等价于 $\mathcal O_X$，当且仅当 $\varphi_n\lambda\in\mathbf{Pic}_{X/k}^0$。因此由定义 9.6.1 和 9.6.8，$\mathcal L$ 与 $\mathcal O_X$ 为 $\tau$-等价，当且仅当 $\lambda\in\mathbf{Pic}_{X/k}^\tau$。\hfill$\square$
\end{origstatement}

\begin{origstatement}{解答}{9.6.13}\label{appA:ans9613}
定理 9.4.8 说明 $\mathbf{Pic}_{X/k}$ 存在并表示 $\Pic_{(X/k)\,(\text{ét})}$。由命题 9.6.12，$\mathbf{Pic}_{X/k}^\tau$ 有限型。因此由习题 9.4.11，$\mathbf{Pic}_{X/k}^\tau$ 拟射影。

再假设 $X$ 几何正规。由于 $\mathbf{Pic}_{X/k}^\tau$ 拟射影，为证明它射影，只需证明它完备。由命题 9.6.12，$\mathbf{Pic}_{X/k}^0$ 的形成与扩张 $k$ 可交换。由 [EGAIV2, 2.7.1(vii)]，一个 $k$-概形完备，当且仅当它在扩张 $k$ 后完备。因此假设 $k$ 代数闭。

当 $\lambda$ 遍历 $\mathbf{Pic}_{X/k}^\tau$ 的 $k$-点时，陪集 $\lambda\mathbf{Pic}_{X/k}^0$ 覆盖 $\mathbf{Pic}_{X/k}^\tau$。由命题 9.5.3，$\mathbf{Pic}_{X/k}^0$ 开，且 $\mathbf{Pic}_{X/k}^\tau$ 拟紧，所以有限多个陪集即可覆盖。由定理 9.5.4，$\mathbf{Pic}_{X/k}^0$ 射影，因而完备。再由命题 9.6.12，$\mathbf{Pic}_{X/k}^\tau$ 闭。因此 $\mathbf{Pic}_{X/k}^\tau$ 完备。\hfill$\square$
\end{origstatement}

\begin{origstatement}{解答}{9.6.15}\label{appA:ans9615}
先假设 $L$ 有界。那么 $\mathcal M$ 定义映射 $\theta\colon T\to\mathbf{Pic}_{X/S}$，且 $\theta(T)\supset\Lambda$。由于 $T$ Noether，显然 $\theta(T)$ 也Noether，因而其任意子空间也Noether。因此 $\Lambda$ 拟紧。

反过来，假设 $\Lambda$ 拟紧。由命题 9.4.17，$\mathbf{Pic}_{X/S}$ 局部有限型，所以 $\Lambda$ 的任意给定点都包含在某个有限型开子概形中。因此有限多个这样的子概形覆盖 $\Lambda$，以 $U$ 表示它们的并。

对某个 fppf 覆盖 $T\to U$，包含 $U\hookrightarrow\mathbf{Pic}_{X/S}$ 由 $X_T$ 上的可逆层 $\mathcal M$ 定义。用 $T$ 的某个开子概形代替 $T$，使 $T\to U$ 有限型且满射。由于 $U$ 有限型，$T$ 也有限型。给定 $\lambda\in\Lambda$，取映到 $\lambda$ 的 $t\in T$，则 $\lambda$ 对应于 $\mathcal M_t$ 的类。\hfill$\square$
\end{origstatement}

\begin{origstatement}{解答}{9.6.18}\label{appA:ans9618}
定理 9.6.16 断言 $\mathbf{Pic}_{X/S}^\tau$ 有限型。因此由习题 9.5.7，它射影。\hfill$\square$
\end{origstatement}

% Source PDF 324 = printed 312.
\begin{origstatement}{解答}{9.6.21}\label{appA:ans9621}
显然，用 $S$ 的开子集代替 $S$，可假设 $X/S$ 射影且 $S$ 连通。给定 $s\in S$，置 $\psi(n)\defeq\chi(\mathcal O_{X_s}(n))$，则 $\psi(n)$ 不依赖于 $s$。给定 $m$，置 $\varphi(n)\defeq m+\psi(n)$。

设 $\lambda\in\mathbf{Pic}_{X/S}$。那么 $\lambda\in\mathbf{Pic}_{X/S}^m$，当且仅当 $\lambda$ 表示次数为 $m$ 的可逆层 $\mathcal L$；而 $\lambda\in\mathbf{Pic}_{X/S}^{\varphi}$，当且仅当 $\chi(\mathcal L(n))=\varphi(n)$。由 Riemann 定理及 $\deg(\bullet)$ 的可加性，
\[
 \chi(\mathcal L(n))=\deg(\mathcal L(n))+\psi(0)=\deg(\mathcal L)+\psi(n).
\]
因此 $\mathbf{Pic}_{X/S}^m=\mathbf{Pic}_{X/S}^{\varphi}$。所以除关于 $\mathbf{Pic}_{X/S}^0$ 的中间两个断言外，其余断言均由定理 9.6.20 得出。

为证明 $\mathbf{Pic}_{X/S}^0=\mathbf{Pic}_{X/S}^{\tau}$，类似地，只需证明 $\deg\mathcal L=0$ 当且仅当 $\mathcal L$ 与 $\mathcal O_X$ 为 $\tau$-等价，因为由习题 9.6.11，后一个条件等价于 $\lambda\in\mathbf{Pic}_{X/S}^{\tau}$。显然可假设 $\mathcal L$ 定义在 $X/S$ 的一个几何纤维上。那么由定理 9.6.3，$\mathcal L$ 的这两个条件等价。

由于 $\deg$ 可加，乘法在集合论意义下把 $\mathbf{Pic}_{X/S}^0\times\mathbf{Pic}_{X/S}^m$ 映入 $\mathbf{Pic}_{X/S}^m$。由于这两个集合在 $\mathbf{Pic}_{X/S}$ 中开，$\mathbf{Pic}_{X/S}^0$ 作用于 $\mathbf{Pic}_{X/S}^m$。

由于 $X/S$ 平坦且几何纤维整，其光滑轨迹 $X_0$ 给出 $S$ 的 fppf 覆盖。暂时作基变换 $X_0\to S$。之后，新的映射 $X_0\to S$ 有截面，其像是相对有效除子 $D$，而与 $\mathcal O_X(mD)$ 作张量积给出所需的从 $\mathbf{Pic}_{X/S}^0$ 到 $\mathbf{Pic}_{X/S}^m$ 的同构。

最后，为说明记号没有歧义，必须证明纤维 $(\mathbf{Pic}_{X/S}^0)_s$ 连通。为此，改作到底域为包含 $k_s$ 的代数闭域 $k$ 的谱的基变换。那么 $X_0$ 有 $k$-有理点 $D$，再次与 $\mathcal O_X(mD)$ 作张量积，给出从 $\mathbf{Pic}_{X/k}^0$ 到 $\mathbf{Pic}_{X/k}^m$ 的同构。因此，只需证明对某个 $m\geq1$，$\mathbf{Pic}_{X/k}^m$ 连通。

令 $\beta\colon X_0^m\to\mathbf{Div}_{X/k}^m\to\mathbf{Pic}_{X/S}^m$ 为习题 9.3.8 中的映射 $\alpha$ 与 Abel 映射的复合。由于 $X$ 整，$m$ 重积 $X_0^m$ 也整。因此，只需证明对某个 $m\geq1$，$\beta$ 满射。

由习题 9.6.7，存在 $m_0\geq1$，使 $X$ 上每个次数为 $0$ 的可逆层都是 $m_0$-正则的。置 $m\defeq\deg(\mathcal O_X(m_0))$。那么 $X$ 上每个次数为 $m$ 的可逆层 $\mathcal L$ 都是 $0$-正则的，因而由整体截面生成。

特别地，对 $X$ 的每个奇点，都有一个整体截面在该点不消失。因此，由于 $k$ 无限，取这些截面的一般线性组合，就在任何奇点都不消失。这个组合定义有效除子 $E$，使 $\mathcal O_X(E)=\mathcal L$。所以 $\beta$ 满射，正如所求。\hfill$\square$
\end{origstatement}

\begin{origstatement}{解答}{9.6.29}\label{appA:ans9629}
假设 $\Lambda$ 拟紧。由习题 9.6.15，存在有限型 $S$-概形 $T$ 及 $X_T$ 上的可逆层 $\mathcal M$，使每个多项式 $\varphi\in\Pi$ 都形如 $\varphi(n)=\chi(\mathcal M_t(n))$，其中 $t\in T$。因此由 [EGAIII2, 7.9.4]，$\varphi$ 的个数至多是 $T$ 的连通分支数。所以 $\Pi$ 有限。

假设 $\Lambda$ 连通，则其闭包也连通，所以可假设 $\Lambda$ 闭。赋予 $\Lambda$ 约化子概形结构。那么对某个 fppf 覆盖 $T\to\Lambda$，包含 $\Lambda\hookrightarrow\mathbf{Pic}_{X/S}$ 由 $X_T$ 上的可逆层 $\mathcal M$ 定义。固定 $t\in T$，置 $\varphi(n)\defeq\chi(\mathcal M_t(n))$。固定 $n$，令 $T'$ 为满足 $\chi(\mathcal M_{t'}(n))=\varphi(n)$ 的点 $t'\in T$ 的集合。由 [EGAIII2, 7.9.4]，$T'$ 开，补集也开。因此它们在 $\Lambda$ 中的像开，且显然互不相交。但 $\Lambda$ 连通，所以 $T'=T$。因此 $\Pi=\{\varphi\}$。\hfill$\square$
\end{origstatement}
% End Appendix A; PDF324 = printed312.
